\subsection{トリガー条件を変えた追加比較}
LTOF各セグメントの左右ANDをORする条件（以下、ペアAND）と、LTOF左右12読み出しのいずれかを許す条件（以下、Flat OR）を比較する。両条件ともUTOF左右のANDを要求する。Flat ORは$(0\land1)\land(2\lor3\lor\cdots\lor13)$とし、0、1がUTOF左右、2--13がLTOF各読み出しである。トリガー幅$30\times4$ ns、LTOF時刻補正$-12\times4$ ns、Slice窓$[-250,+250]\times4$ ns、入力先頭20,000 TF、5段各1プロセス、FileSinkの出力先/dev/null、各条件1回とした。全4台で測定の前後にperformanceガバナーを確認し、HLFバイナリーは先のペアAND測定と同一である。過去発表時のトリガー設定との一致は未確認である。

トリガーレートとSlice数の比較はmarq07のみで行った。これは取得データの時間に対するレートであり、リプレイヤーの再生速度や処理スループットとは異なる。対象TFは連続する4 HBFからなり、HBF周期524.288 $\mu$sから1 TFは2.097152 ms、20,000 TFは41.94304 sである。平均トリガーレートは総トリガー数/41.94304 s、平均Slice数は総Slice数/20,000 TFであり、Sliceが0のTFも分母に含める。

LogicFilterの正のTotalHitsログの全件列と、TF順に数えたHLF入力Slice数の正の列が両run・両条件で一致することを確認した。Slicerはトリガーごとに1 Sliceを生成し、隣接窓が重なる場合も統合せずヒットを再使用する。このため、以下の「トリガー数」はLogicFilter出力の数であり、独立な物理事象数を意味しない。0 Slice TFを独立なスピルoff判定と同一視しない。

時系列は取得開始HBFからの相対時間で0.1 sごとに集計する。各TFの全トリガー数をTF中央時刻のビンに割り当てるため、ビン境界で最大約1 TFの時刻割当誤差を含む。各トリガーの時刻を直接0.1 s境界で選別した図ではない。最後の0.04304 sもその実時間で割って時系列に示し、0.1 sレート分布からは除く。ヒストグラムは200ビン、軸は線形、グリッドを使用した。
\begin{table}[htbp]\centering\small
\caption{トリガー条件による平均値の比較（marq07）。いずれも全20,000 TF、41.94304 sを対象とする。}
\begin{tabular}{rlrrrr}\hline
run & LTOF条件 & 総トリガー数 & 平均Slice/TF & 平均レート [kHz] & 0 Slice TF \\\hline
400 & ペアAND & 156,792 & 7.8396 & 3.738 & 10,221 \\
400 & Flat OR & 2,957,669 & 147.8835 & 70.516 & 9,258 \\
668 & ペアAND & 151,862 & 7.5931 & 3.621 & 10,340 \\
668 & Flat OR & 2,919,747 & 145.9873 & 69.612 & 9,456 \\
\hline\end{tabular}\label{tab:trigger-logic-comparison}\end{table}
\input{\subfix{../../figure/05-hlf/flat-or-run000668-summary/tables.tex}}

\begin{figure}[htbp]
\centering\makebox[\linewidth][c]{\includegraphics[pagebox=cropbox,width=.90\linewidth,height=.65\textheight,keepaspectratio]{05-hlf/trigger-logic-run000668/trigger-time-series.pdf}}
\caption{run668のトリガーレート推移（marq07）。青：LTOFペアAND後OR、赤：LTOF個別読み出しのFlat OR。0.1 s集計、最終点のみ0.04304 s。}
\label{fig:trigger-logic-668-time}
\end{figure}

\begin{figure}[htbp]
\centering\makebox[\linewidth][c]{\includegraphics[pagebox=cropbox,width=.46\linewidth,height=.65\textheight,keepaspectratio]{05-hlf/trigger-logic-run000668/slice-count.pdf}\includegraphics[pagebox=cropbox,width=.46\linewidth,height=.65\textheight,keepaspectratio]{05-hlf/trigger-logic-run000668/tf-trigger-rate.pdf}}
\caption{run668のHLF入力Slice数/TF（左）とTF時間換算のトリガーレート（右）。全TFを含む、200ビン。}
\label{fig:trigger-logic-668-slices}
\end{figure}

\begin{figure}[htbp]
\centering\makebox[\linewidth][c]{\includegraphics[pagebox=cropbox,width=.90\linewidth,height=.65\textheight,keepaspectratio]{05-hlf/trigger-logic-run000668/100ms-trigger-rate.pdf}}
\caption{run668の0.1 s平均トリガーレート分布。419個の完全な0.1 s窓、200ビン。}
\label{fig:trigger-logic-668-100ms}
\end{figure}

\begin{figure}[htbp]
\centering\makebox[\linewidth][c]{\includegraphics[pagebox=cropbox,width=.46\linewidth,height=.65\textheight,keepaspectratio]{05-hlf/flat-or-run000668-input-rate/raw-tf-size.pdf}\includegraphics[pagebox=cropbox,width=.46\linewidth,height=.65\textheight,keepaspectratio]{05-hlf/flat-or-run000668-input-rate/hlf-tf-size.pdf}}
\caption{Flat OR条件、run000668の1 TFのデータサイズ。左：LF前生データ、右：HLF入力。Slice有無はHLF入力での分類。}
\label{fig:flat-668-tf-size}
\end{figure}

\begin{figure}[htbp]
\centering\makebox[\linewidth][c]{\includegraphics[pagebox=cropbox,width=.46\linewidth,height=.65\textheight,keepaspectratio]{05-hlf/flat-or-run000668-input-rate/raw-tf-rate.pdf}\includegraphics[pagebox=cropbox,width=.46\linewidth,height=.65\textheight,keepaspectratio]{05-hlf/flat-or-run000668-input-rate/hlf-tf-rate.pdf}}
\caption{Flat OR条件、run000668のTF時間で換算した瞬間取得レート。左：LF前生データ、右：HLF入力。Slice有無はHLF入力での分類。}
\label{fig:flat-668-tf-rate}
\end{figure}

\begin{figure}[htbp]
\centering\makebox[\linewidth][c]{\includegraphics[pagebox=cropbox,width=.46\linewidth,height=.65\textheight,keepaspectratio]{05-hlf/flat-or-run000668-input-rate/raw-100ms-rate.pdf}\includegraphics[pagebox=cropbox,width=.46\linewidth,height=.65\textheight,keepaspectratio]{05-hlf/flat-or-run000668-input-rate/hlf-100ms-rate.pdf}}
\caption{Flat OR条件、run000668の0.1 sごとの取得レート。左：LF前生データ、右：HLF入力。Slice有無はHLF入力での分類。}
\label{fig:flat-668-100ms-rate}
\end{figure}

\begin{figure}[htbp]
\centering\makebox[\linewidth][c]{\includegraphics[pagebox=cropbox,width=.46\linewidth,height=.65\textheight,keepaspectratio]{05-hlf/flat-or-run000668-tof/cond-all-size-time-e50server03.pdf}\includegraphics[pagebox=cropbox,width=.46\linewidth,height=.65\textheight,keepaspectratio]{05-hlf/flat-or-run000668-tof/slice-has_slice-size-time-e50server03.pdf}}
\caption{Flat OR条件、run000668、tof、e50server03。左：ConditionalRun全TF、右：ProcessSlice。横軸：入力サイズ、縦軸：処理時間。200×200ビン、線形軸、グリッド。}
\label{fig:flat-668-tof-e50server03}
\end{figure}

\begin{figure}[htbp]
\centering\makebox[\linewidth][c]{\includegraphics[pagebox=cropbox,width=.46\linewidth,height=.65\textheight,keepaspectratio]{05-hlf/flat-or-run000668-tof/cond-all-size-time-marq04.pdf}\includegraphics[pagebox=cropbox,width=.46\linewidth,height=.65\textheight,keepaspectratio]{05-hlf/flat-or-run000668-tof/slice-has_slice-size-time-marq04.pdf}}
\caption{Flat OR条件、run000668、tof、marq04。左：ConditionalRun全TF、右：ProcessSlice。横軸：入力サイズ、縦軸：処理時間。200×200ビン、線形軸、グリッド。}
\label{fig:flat-668-tof-marq04}
\end{figure}

\begin{figure}[htbp]
\centering\makebox[\linewidth][c]{\includegraphics[pagebox=cropbox,width=.46\linewidth,height=.65\textheight,keepaspectratio]{05-hlf/flat-or-run000668-tof/cond-all-size-time-marq06.pdf}\includegraphics[pagebox=cropbox,width=.46\linewidth,height=.65\textheight,keepaspectratio]{05-hlf/flat-or-run000668-tof/slice-has_slice-size-time-marq06.pdf}}
\caption{Flat OR条件、run000668、tof、marq06。左：ConditionalRun全TF、右：ProcessSlice。横軸：入力サイズ、縦軸：処理時間。200×200ビン、線形軸、グリッド。}
\label{fig:flat-668-tof-marq06}
\end{figure}

\begin{figure}[htbp]
\centering\makebox[\linewidth][c]{\includegraphics[pagebox=cropbox,width=.46\linewidth,height=.65\textheight,keepaspectratio]{05-hlf/flat-or-run000668-tof/cond-all-size-time-marq07.pdf}\includegraphics[pagebox=cropbox,width=.46\linewidth,height=.65\textheight,keepaspectratio]{05-hlf/flat-or-run000668-tof/slice-has_slice-size-time-marq07.pdf}}
\caption{Flat OR条件、run000668、tof、marq07。左：ConditionalRun全TF、右：ProcessSlice。横軸：入力サイズ、縦軸：処理時間。200×200ビン、線形軸、グリッド。}
\label{fig:flat-668-tof-marq07}
\end{figure}

\begin{figure}[htbp]
\centering\makebox[\linewidth][c]{\includegraphics[pagebox=cropbox,width=.46\linewidth,height=.65\textheight,keepaspectratio]{05-hlf/flat-or-run000668-multiplicity/cond-all-size-time-e50server03.pdf}\includegraphics[pagebox=cropbox,width=.46\linewidth,height=.65\textheight,keepaspectratio]{05-hlf/flat-or-run000668-multiplicity/slice-has_slice-size-time-e50server03.pdf}}
\caption{Flat OR条件、run000668、multiplicity、e50server03。左：ConditionalRun全TF、右：ProcessSlice。横軸：入力サイズ、縦軸：処理時間。200×200ビン、線形軸、グリッド。}
\label{fig:flat-668-multiplicity-e50server03}
\end{figure}

\begin{figure}[htbp]
\centering\makebox[\linewidth][c]{\includegraphics[pagebox=cropbox,width=.46\linewidth,height=.65\textheight,keepaspectratio]{05-hlf/flat-or-run000668-multiplicity/cond-all-size-time-marq04.pdf}\includegraphics[pagebox=cropbox,width=.46\linewidth,height=.65\textheight,keepaspectratio]{05-hlf/flat-or-run000668-multiplicity/slice-has_slice-size-time-marq04.pdf}}
\caption{Flat OR条件、run000668、multiplicity、marq04。左：ConditionalRun全TF、右：ProcessSlice。横軸：入力サイズ、縦軸：処理時間。200×200ビン、線形軸、グリッド。}
\label{fig:flat-668-multiplicity-marq04}
\end{figure}

\begin{figure}[htbp]
\centering\makebox[\linewidth][c]{\includegraphics[pagebox=cropbox,width=.46\linewidth,height=.65\textheight,keepaspectratio]{05-hlf/flat-or-run000668-multiplicity/cond-all-size-time-marq06.pdf}\includegraphics[pagebox=cropbox,width=.46\linewidth,height=.65\textheight,keepaspectratio]{05-hlf/flat-or-run000668-multiplicity/slice-has_slice-size-time-marq06.pdf}}
\caption{Flat OR条件、run000668、multiplicity、marq06。左：ConditionalRun全TF、右：ProcessSlice。横軸：入力サイズ、縦軸：処理時間。200×200ビン、線形軸、グリッド。}
\label{fig:flat-668-multiplicity-marq06}
\end{figure}

\begin{figure}[htbp]
\centering\makebox[\linewidth][c]{\includegraphics[pagebox=cropbox,width=.46\linewidth,height=.65\textheight,keepaspectratio]{05-hlf/flat-or-run000668-multiplicity/cond-all-size-time-marq07.pdf}\includegraphics[pagebox=cropbox,width=.46\linewidth,height=.65\textheight,keepaspectratio]{05-hlf/flat-or-run000668-multiplicity/slice-has_slice-size-time-marq07.pdf}}
\caption{Flat OR条件、run000668、multiplicity、marq07。左：ConditionalRun全TF、右：ProcessSlice。横軸：入力サイズ、縦軸：処理時間。200×200ビン、線形軸、グリッド。}
\label{fig:flat-668-multiplicity-marq07}
\end{figure}

\begin{figure}[htbp]
\centering\makebox[\linewidth][c]{\includegraphics[pagebox=cropbox,width=.46\linewidth,height=.65\textheight,keepaspectratio]{05-hlf/flat-or-run000668-track/cond-all-size-time-e50server03.pdf}\includegraphics[pagebox=cropbox,width=.46\linewidth,height=.65\textheight,keepaspectratio]{05-hlf/flat-or-run000668-track/slice-has_slice-size-time-e50server03.pdf}}
\caption{Flat OR条件、run000668、track、e50server03。左：ConditionalRun全TF、右：ProcessSlice。横軸：入力サイズ、縦軸：処理時間。200×200ビン、線形軸、グリッド。}
\label{fig:flat-668-track-e50server03}
\end{figure}

\begin{figure}[htbp]
\centering\makebox[\linewidth][c]{\includegraphics[pagebox=cropbox,width=.46\linewidth,height=.65\textheight,keepaspectratio]{05-hlf/flat-or-run000668-track/cond-all-size-time-marq04.pdf}\includegraphics[pagebox=cropbox,width=.46\linewidth,height=.65\textheight,keepaspectratio]{05-hlf/flat-or-run000668-track/slice-has_slice-size-time-marq04.pdf}}
\caption{Flat OR条件、run000668、track、marq04。左：ConditionalRun全TF、右：ProcessSlice。横軸：入力サイズ、縦軸：処理時間。200×200ビン、線形軸、グリッド。}
\label{fig:flat-668-track-marq04}
\end{figure}

\begin{figure}[htbp]
\centering\makebox[\linewidth][c]{\includegraphics[pagebox=cropbox,width=.46\linewidth,height=.65\textheight,keepaspectratio]{05-hlf/flat-or-run000668-track/cond-all-size-time-marq06.pdf}\includegraphics[pagebox=cropbox,width=.46\linewidth,height=.65\textheight,keepaspectratio]{05-hlf/flat-or-run000668-track/slice-has_slice-size-time-marq06.pdf}}
\caption{Flat OR条件、run000668、track、marq06。左：ConditionalRun全TF、右：ProcessSlice。横軸：入力サイズ、縦軸：処理時間。200×200ビン、線形軸、グリッド。}
\label{fig:flat-668-track-marq06}
\end{figure}

\begin{figure}[htbp]
\centering\makebox[\linewidth][c]{\includegraphics[pagebox=cropbox,width=.46\linewidth,height=.65\textheight,keepaspectratio]{05-hlf/flat-or-run000668-track/cond-all-size-time-marq07.pdf}\includegraphics[pagebox=cropbox,width=.46\linewidth,height=.65\textheight,keepaspectratio]{05-hlf/flat-or-run000668-track/slice-has_slice-size-time-marq07.pdf}}
\caption{Flat OR条件、run000668、track、marq07。左：ConditionalRun全TF、右：ProcessSlice。横軸：入力サイズ、縦軸：処理時間。200×200ビン、線形軸、グリッド。}
\label{fig:flat-668-track-marq07}
\end{figure}
e50server03のTOF ConditionalRun全TFの合算スループットは、run000400でペアANDの390.3 MB/sからFlat ORの621.5 MB/s、run000668で360.4 MB/sから606.6 MB/sに変化した。トリガー条件によりHLFへ渡すSlice数と入力サイズが変わるため、これは同一入力サイズでのCPU処理能力の向上を意味しない。旧発表値594.9 MB/sに近い値となったが、当時の設定・計時区間・集計法との一致は未確認であり、再現できたとは断定しない。

0.1 s取得データレートの図では、生データをHBF境界まで用いて配分し、HLF入力サイズはTF内の4 HBFへ等分する近似を用いた。トリガーレートの図はTF中央時刻への割当であり、このサイズ配分の図とは集計手順が異なる。
